\section*{Chapter \ref{ch:s2:government-bond-relative-value} --- Government-Bond Relative Value}

\begin{solution}{exo:s2:government-bond-relative-value:1}
$\ln 0.5/\ln 0.5^{1/40} = 40$ days.
\end{solution}

\begin{solution}{exo:s2:government-bond-relative-value:2}
$2 \times 4.99 - 4.85 - 5.11 = 0.02$ percentage points: 2.0 basis points.
\end{solution}

\begin{solution}{exo:s2:government-bond-relative-value:3}
Least squares makes the residuals orthogonal to the regressors: $X^\top(y - X\hat\beta) = 0$. A book with weights proportional to the residuals is
already neutral to the fitted factors; the hedge matters once the book selects only some bonds.
\end{solution}

\begin{solution}{exo:s2:government-bond-relative-value:4}
The noise adds variance to today's residual without any persistence, so the regression of tomorrow on today is biased towards zero: the fitted
persistence is too low and the half-life too short.
\end{solution}

\begin{solution}{exo:s2:government-bond-relative-value:5}
They often trade at a premium for their liquidity and their specialness in repo, which would distort a curve for ordinary bonds. The premium
suggests selling the on-the-run issue against a nearby off-the-run one, if its repo cost can be carried.
\end{solution}

\begin{solution}{exo:s2:government-bond-relative-value:6}
Removing three factors takes out almost all curve risk, so gross volatility falls most; but the hedge adds small positions in all forty bonds that
change every week, and at 0.25 basis points the extra turnover costs more than the risk it removes is worth.
\end{solution}

\begin{solution}{exo:s2:government-bond-relative-value:7}
The DV01-neutral book rises from 1.13 to 2.54 and the factor-neutral from 0.69 to 3.48: with cheap execution the better hedge wins. The trade belongs to
whoever trades most cheaply, often a dealer or a fund with low execution and financing costs.
\end{solution}

\begin{solution}{exo:s2:government-bond-relative-value:8}
The 18 days come from a one-day regression biased by quote noise and by the fast part of the error; at 120 days 40\% of a residual is still there.
Closing after a month leaves most of the expected reversion uncollected and pays turnover again.
\end{solution}

\begin{solution}{pb:s2:government-bond-relative-value:1}
\begin{enumerate}
\item The Board's daily Treasury curve from 1961, fitted with Svensson's extension of Nelson--Siegel.
\item Callable and flower bonds, securities within three months of maturity, bills, twenty-year bonds from 1996, and on-the-run and first
off-the-run issues: they carry features, segmentation or liquidity premiums that ordinary bonds do not.
\item 87.9\%, 8.2\% and 2.0\%.
\item Market yield minus the fitted curve's yield at the bond's maturity.
\item Middle against wings, neutral in DV01 and slope; mean $-1.6$ basis points, standard deviation 26.6, half-life 105 days.
\item A ranking of bonds by residual to choose which to buy and sell.
\item 2 basis points reverting in 40 days, 2 in 250 days, seen through 0.5 basis points of noise.
\item Five cheapest long, five richest short, weekly; unhedged (long only), DV01-neutral, or neutral in three factors.
\item 0.36, 3.46 and 5.41 gross; 0.25, 1.13 and 0.69 net.
\item It is mostly a duration position.
\item The factor hedge's turnover costs more at 0.25 basis points than its risk reduction is worth; at 0.1 it does not.
\item 81\% left after 20 days, 59\% after 60, 40\% after 120.
\item \textbf{Named result.} Net of costs the DV01-neutral book earns a Sharpe ratio of 1.13 and the factor-neutral book 0.69 (5.41 gross); the
residuals' one-day half-life is 17.9 days, but 40\% of a residual remains after 120 days.
\item Quote noise and the error's slow part both escape it.
\item Financing: haircuts, repo rolls, specialness, crowded positions.
\item Bid-ask, repo and the hedge's turnover.
\item Fit on quotes known then, trade at bid and offer, include repo, measure decay at the holding horizon.
\item Off-the-run versus on-the-run.
\item Chapter~11 trades a bond against its future, financed in repo at high leverage.
\item Liquidity and balance sheet to investors whose demands push bonds away from the curve.
\end{enumerate}
\end{solution}

\begin{solution}{iq:s2:government-bond-relative-value:1}
Its yield against a curve fitted to comparable bonds: supply, index and regulatory demand, liquidity, repo specialness, coupon and tax effects make
bonds sit above or below.
\end{solution}

\begin{solution}{iq:s2:government-bond-relative-value:2}
A smooth parametric form (Nelson--Siegel--Svensson) or a spline fitted to seasoned notes and bonds weighted by duration; leave out bills, bonds near
maturity, on-the-run issues, and bonds with special features.
\end{solution}

\begin{solution}{iq:s2:government-bond-relative-value:3}
Solve for wing weights so that total DV01 and exposure to a slope move are zero; the position still bets on curvature at the middle, on the three
bonds' own residuals, and on their repo rates.
\end{solution}

\begin{solution}{iq:s2:government-bond-relative-value:4}
Haircuts rise and force selling, repo cannot be rolled, residuals widen before they revert, many funds hold the same positions and sell together.
\end{solution}

\begin{solution}{iq:s2:government-bond-relative-value:5}
Inputs: clean prices or yields of eligible bonds; exclusions by rule; a weighted least-squares fit; outputs: parameters, fitted yields, residuals
and their history; checks: fit error, parameter jumps, stale prices, new issues.
\end{solution}

\begin{solution}{iq:s2:government-bond-relative-value:6}
The slope is $\mathrm{Cov}(e_{t+1} + \eta_{t+1}, e_t + \eta_t)/\mathrm{Var}(e_t + \eta_t) = \phi\sigma_e^2/(\sigma_e^2 + \sigma_\eta^2)$. It is below
$\phi$, so $\ln 0.5/\ln(\text{slope})$ understates the half-life; the bias grows with the noise's share of the residual's variance.
\end{solution}
